\clearpage\fancyhead[L]{\small PALLAS.}\hypertarget{native-361-lower}{}\label{native-361-lower}
\begin{center}\textbf{PALLAS.}\par\vspace{4pt}\hrule\vspace{4pt}\fontsize{10}{12}\selectfont Göttingische gelehrte Anzeigen. Heft 127. S. 1257\;..\;1259. 1812 Aug. 8.\par\vspace{4pt}\hrule\vspace{4pt}\end{center}

\par\indent In Heft 67 dieser Zeitschrift [S. 349 dieses Bandes] wurde bereits darüber berichtet, dass Pallas in diesem Jahr zum ersten Mal wiedergefunden wurde und ihre Stellung sehr gut mit der nach der Störungstheorie berechneten Ephemeride übereinstimmte. Seither wurden an der hiesigen Sternwarte nur wenige Beobachtungen gemacht. Vor allem um die Zeit der Opposition ließen ungünstiges Wetter, Mondschein und der hohe, für unseren Standort ungünstige Stand des Planeten nur eine einzige Beobachtung zu. Hier folgen sämtliche Beobachtungen, die Professor \textsc{Gauss} in diesem Jahr an unserer Sternwarte gemacht hat:

\smallskip{\centering
\fontsize{10}{12}\selectfont\setlength{\tabcolsep}{3pt}\renewcommand{\arraystretch}{1.15}\begin{tabular}{c|c|c|c}
1812 & Mittlere Zeit & \shortstack{Rektaszension} & \shortstack{Nördliche\\Abweichung} \\
\hline
April 9 & $11^{\mathrm u}\;32^{\mathrm m}\;27^{\mathrm s}$ & $269^\circ\;5'\;58''4$ & $16^\circ\;38'\;0''$ \\
Mai 2 & 10 \quad 36 \quad 6 & 268 \quad 42 \quad 5.3 & 21 \quad 4 \quad 31.7 \\
— 4 & 10 \quad 44 \quad 44 & 268 \quad 31 \quad 38.0 & 21 \quad 25 \quad 4.3 \\
Juni 7 & 11 \quad 13 \quad 51 & 262 \quad 50 \quad 51.1 & 25 \quad 2 \quad 56.6 \\
\end{tabular}\par}
\smallskip

\par\indent Bei der letzten Beobachtung hatte der Planet die zehnte Größe.

\par\indent Professor \textsc{Gauss} hatte mehrere seiner auswärtigen Astronomenfreunde, denen bessere Hilfsmittel zur Verfügung stehen, eingeladen, Pallas um die Zeit der Opposition zu beobachten. Nach den bisher eingegangenen Nachrichten ist dies jedoch weder an der kaiserlichen Sternwarte in Paris noch in Marseille oder Berlin gelungen. Das ist kaum verwunderlich, denn der Planet ist dieses Jahr so lichtschwach, dass er an den fest montierten Instrumenten kaum die geringste Beleuchtung verträgt. Nur Herr \textsc{Burckhardt} hatte das Glück, an der Sternwarte der Militärschule in Paris die folgenden Meridianbeobachtungen zu erhalten:

\smallskip{\centering
\fontsize{10}{12}\selectfont\setlength{\tabcolsep}{3pt}\renewcommand{\arraystretch}{1.15}\begin{tabular}{c|c|c|c}
1812 & Mittlere Zeit & \shortstack{Rektaszension} & \shortstack{Nördliche\\Abweichung} \\
\hline
Juni 5 & $12^{\mathrm u}\;35^{\mathrm m}\;57^{\mathrm s}4$ & $263^\circ\;15'\;24''0$ & $24^\circ\;58'\;28''5$ \\
6 & 12 \quad 33 \quad 10.5 & 263 \quad 2 \quad 35 & 25 \quad 0 \quad 30 \\
8 & 12 \quad 21 \quad 37.0 & 262 \quad 37 \quad 6.5 & 25 \quad 3 \quad 57.7 \\
11 & 12 \quad 7 \quad 15.7 & 261 \quad 58 \quad 37.8 & 25 \quad 6 \quad 29.3 \\
\end{tabular}\par}
\smallskip

\par\indent Professor \textsc{Gauss} übergab diese Beobachtungen Herrn \textsc{Nicolai}, damit er daraus die Opposition berechne. Sein Ergebnis lautet:
\par\smallskip\noindent\SourceLeaf{362}\begin{minipage}{\linewidth}\centering
\fontsize{10}{12}\selectfont\setlength{\tabcolsep}{0pt}\renewcommand{\arraystretch}{1.15}\begin{tabular}{l}
Opposition der Pallas 1812 Juni 10. $3^{\mathrm u}\;2^{\mathrm m}\;48^{\mathrm s}$\\mittl. Zeit am Göttinger Meridian
\end{tabular}\par\smallskip
\fontsize{10}{12}\selectfont\setlength{\tabcolsep}{2pt}\renewcommand{\arraystretch}{1.15}\begin{tabular}{lcr}
Wahre Länge & $.\;.\;.\;.\;.\;.\;.\;.\;.\;.\;.\;.\;.\;.\;.\;.\;.\;.$ & $259^\circ\;37'\;51''7$\\
Wahre geozentrische Breite & $.\;.\;.\;.\;.\;.\;.\;.\;.\;.\;.\;.\;.$ & $48\;16\;26.1\;\mathrm{N}.$\\
\end{tabular}\par\end{minipage}\par

\par\indent Dies ist nun die \emph{achte} beobachtete Opposition der Pallas seit ihrer Entdeckung. Das Ergebnis weicht von der Rechnung, die Herr \textsc{Nicolai} ein Jahr zuvor im Voraus durchgeführt hat, um 29 Minuten in der Zeit, 1 Minute in der Länge und 10 Sekunden in der Breite ab (siehe Gött. gel. Anz. 1811, S. 1292 [S. 329 dieses Bandes]). Professor \textsc{Gauss} glaubt, in diesen zwar geringen Unterschieden schon ganz sicher den Einfluss der bisher noch nicht untersuchten Störungen durch Saturn und Mars zu erkennen. Für die Berechnung der Störungen durch Saturn sind bereits alle Vorbereitungen getroffen, und wir werden darüber zu gegebener Zeit ausführlich berichten. Inzwischen hat Professor \textsc{Gauss} untersucht, wie genau sich alle acht bisher beobachteten Oppositionen noch darstellen lassen, wenn nur die Störungen durch Jupiter berücksichtigt werden. Das Ergebnis ist, dass man für die größtmögliche Übereinstimmung

\begin{center}\fontsize{9}{11}\selectfont\setlength{\tabcolsep}{2pt}\renewcommand{\arraystretch}{1.15}\begin{tabular}{lcr}
\shortstack[l]{die Epoche der mittleren Länge\\am 10. Juni 1812 um} & $.\;.\;.\;.$ & $-26''95$\\
die tägliche mittlere Bewegung um & $.\;.\;.\;.\;.\;.\;.\;.\;.\;.$ & $-0''01225$\\
den Logarithmus der Halbachse um & $.\;.\;.\;.\;.\;.\;.\;.\;.\;.\;.$ & $+0.0000046$\\
die Länge des Perihels um & $.\;.\;.\;.\;.\;.\;.\;.\;.\;.\;.\;.\;.\;.$ & $-18''84$\\
den Exzentrizitätswinkel um & $.\;.\;.\;.\;.\;.\;.\;.\;.\;.\;.\;.\;.$ & $-3''92$\\
die Länge des Knotens um & $.\;.\;.\;.\;.\;.\;.\;.\;.\;.\;.\;.\;.\;.$ & $+0''80$\\
die Neigung der Bahn um & $.\;.\;.\;.\;.\;.\;.\;.\;.\;.\;.\;.\;.\;.$ & $-1''21$\\
\end{tabular}\end{center}

\Needspace{16\baselineskip}\noindent ändern muss. Die beobachtete und die berechnete Opposition stimmen dann wie folgt überein:

\begin{center}
\fontsize{10}{12}\selectfont\setlength{\tabcolsep}{7pt}\renewcommand{\arraystretch}{1.15}\begin{tabular}{c|c|c}
 & \multicolumn{2}{c}{Unterschied im} \\
Opposition & mittlere Länge & heliozentr. Breite \\
\hline
1803 & $+$\ 6$''$4  & $-$\ 2$''$8 \\
1804 & $-$10.9 & $-$9.7 \\
1805 & $-$13.8 & $-$0.7 \\
1807 & $+$10.2 & $-$5.8 \\
1808 & $+$19.1 & $+$11.7 \\
1809 & $+$16.1 & $-$1.5 \\
1811 & $+$14.2 & $-$9.5 \\
1812 & $-$22.3 & $+$6.1 \\
\end{tabular}
\end{center}

\par\indent Die \emph{neunte} Opposition wird nach der Rechnung von Herrn \textsc{Nicolai} am 18. August 1813 um $9^{\mathrm u}$ eintreten, bei $325^\circ\;24'$ Länge und $24^\circ\;37'$ nördlicher geozentrischer Breite.
